\documentclass[11pt]{article}
\usepackage[margin=1in]{geometry}
\usepackage{amsmath,amssymb,amsthm}
\usepackage[hidelinks]{hyperref}
\usepackage{microtype}

\title{\bfseries Orbital Dragons to Hunt Cryptids:\\
Neural Networks, Quantum-Inspired Data Science,\\
and the Busy Beaver Problem}
\author{Enan Srivastava}
\date{DeepMath 2026 One-Page Abstract}

\newtheorem{theorem}{Theorem}
\newcommand{\Halts}{\mathsf{RawHalts}}
\newcommand{\Safe}{\mathsf{SafeMetaOrbit}}

\begin{document}
\maketitle

\noindent
The busy beaver value $BB(n)$ is the longest finite runtime of any
$n$-state Turing machine started on a blank tape. $BB(5)$ was settled
in 2024 with a computer-checked proof. $BB(6)$ is blocked by a handful
of machines whose tapes secretly run Collatz-style arithmetic; nobody
can prove whether they stop. This paper studies two of these
machines and surveys the rest. The approach is empirical first: the
machine's behavior is logged as a
sequence of events, the event data embedded in logarithmic, simplex, and
hyperbolic coordinates, and the dynamics found to be piecewise
affine in a five-integer state. The description is then proven
correct in the Lean proof assistant, down to the machine's individual steps and
with a minimal axiom footprint. The result is a machine-checked
theorem: if a simple five-integer recurrence never enters an explicit
finite set of bad configurations, the machine never halts. The
converse pressure is also proven: entering a bad configuration forces the
machine to halt. The open problem is thereby moved, without loss, from
a Turing machine to a short arithmetic recurrence. The renewal map of the recurrence is then derived in closed form and shown
to have a width-five fatal window around the diagonal, and the two mechanical
proof routes with computational evidence: no congruence invariant up to
modulus 384 excludes the window, and the branch dynamics has
exponential factor complexity, so it is not conjugate to a rotation.
Measuring the recurrence over its whole seed space shows the set of
dying seeds has natural density zero at the Borel--Cantelli rate
($\approx 12/X$), so the studied machines' seeds are typical
survivors.
Surveying the entire holdout set with a two-ended harness (forward
orbit against backward halt cone, regimes fitted in between)
classifies all $1{,}094$ open six-state machines: zero halts, zero
alignments; four fixed-dimension machines, exactly the certified
ones; an eleven-machine \emph{array class} whose state spaces grow in
dimension, for whose leading member an exact word-rewrite calculus
is constructed (a two-phase odometer with a block-trisection boundary rule)
validated against $2 \times 10^{8}$ raw steps; and $307$ single-law
machines among $346$ with regime structure. The program is framed
on a verification spectrum whose proven endpoint is
$\mathsf{MIP}^{*} = \mathsf{RE}$: kernel checking at the bottom,
certified predictive emulation (this paper) in the middle, and
entangled-prover verification of halting-type claims as the goal. Every
empirical claim is validated against $10^{9}$ raw machine steps, and
the recurrence is tested for $10^{8}$ iterations. For the theory of
learning, the pipeline is representation learning carried to a
proof, and it is offered as a benchmark: unbounded, exactly labeled,
law-governed sequence data on which a learned representation has a
provable ground truth, with a two-prover game that verifies the
forecasts of any learned emulator at constant cost per check.

\medskip
\noindent The central certified statement, one per machine:

\begin{theorem}
Checked by the Lean kernel at axioms
$\{\mathrm{propext},\mathrm{Quot.sound}\}$:
$\ \Safe \rightarrow \lnot\Halts(\mathsf{blankStart})$,
and configurations outside the guarded region halt within about
fifty steps, so $\Safe$ is the halting problem restated as
arithmetic.
\end{theorem}

\newpage

\begin{thebibliography}{99}\small
\bibitem{coqbb5} The bbchallenge Collaboration.
Determination of the fifth Busy Beaver value. arXiv:2509.12337.
\bibitem{deciders} The bbchallenge Collaboration.
Turing machine deciders, part I. arXiv:2504.20563.
\bibitem{michel} P.~Michel. The Busy Beaver Competition: a
historical survey. arXiv:0906.3749.
\bibitem{aaronson} S.~Aaronson. The Busy Beaver Frontier.
\emph{SIGACT News} 51(3), 2020.
\bibitem{tao-collatz} T.~Tao. Almost all orbits of the Collatz map
attain almost bounded values. arXiv:1909.03562.
\bibitem{moore} C.~Moore. Generalized shifts: unpredictability and
undecidability in dynamical systems. \emph{Nonlinearity} 4 (1991).
\bibitem{kurka} P.~K\r{u}rka. On topological dynamics of Turing
machines. \emph{Theoret.\ Comput.\ Sci.} 174 (1997).
\bibitem{viro} O.~Viro. Dequantization of real algebraic geometry on
logarithmic paper. arXiv:math/0005163.
\bibitem{litvinov} G.~Litvinov. The Maslov dequantization, idempotent
and tropical mathematics. arXiv:math/0507014.
\bibitem{nickel-kiela} M.~Nickel and D.~Kiela. Poincar\'e embeddings
for learning hierarchical representations. NeurIPS 2017.
\bibitem{series} C.~Series. The modular surface and continued
fractions. \emph{J.\ London Math.\ Soc.} 31 (1985).
\bibitem{mipre} Z.~Ji, A.~Natarajan, T.~Vidick, J.~Wright, and
H.~Yuen. $\mathsf{MIP}^{*} = \mathsf{RE}$. arXiv:2001.04383;
\emph{Comm.\ ACM} 64(11), 2021.
\bibitem{mnyu} H.~Mousavi, S.~Nezhadi, and H.~Yuen. Nonlocal games,
compression theorems, and the arithmetical hierarchy. STOC 2022;
arXiv:2110.04651.
\bibitem{etp} T.~Tao et al. The Equational Theories Project.
arXiv:2512.07087.
\bibitem{palomar} Palomar: a registry of Lean-verified
mathematics, 2026. \texttt{palomar-registry.org}.
\end{thebibliography}

\end{document}
